\section{Linear deep-water waves and dispersion}
\label{sec:dispersion}
\subsection{Deriving the carrier dispersion relation}
Linearize about a flat, resting surface and evaluate at $z=0$:
\begin{equation}\label{eq:linear-bc}
 \zeta_t=\phi_z,\qquad \phi_t=-g\zeta.
\end{equation}
Take one complex Fourier component
\[
 \phi=Ce^{k_0z}e^{i(k_0x-\omega_0t)},\qquad
 \zeta=De^{i(k_0x-\omega_0t)}.
\]
Real fields include the conjugates. Laplace's equation holds because
$\partial_x^2$ contributes $-k_0^2$ and $\partial_z^2$ contributes $+k_0^2$.
$e^{k_0z}$ decays as $z\to-\infty$; the growing exponential is excluded.

The surface conditions give
\[
 -i\omega_0D=k_0C,\qquad -i\omega_0C=-gD.
\]
Elimination yields
\begin{equation}\label{eq:carrier-dispersion}
 \boxed{\omega_0^2=gk_0.}
\end{equation}
This also verifies \eqref{eq:AB}. For two horizontal directions, replace
$k_0$ by the horizontal wavevector's magnitude.

\subsection{Phase velocity and group velocity}
A crest has fixed phase $k_0x-\omega_0t$, so $c_p=\omega_0/k_0$.
To find packet velocity, add two nearby waves:
\begin{align*}
 &\cos((k_0+q)x-(\omega_0+\delta\omega)t)
 +\cos((k_0-q)x-(\omega_0-\delta\omega)t)\\
 &\hspace{1cm}=2\cos(qx-\delta\omega t)\cos(k_0x-\omega_0t).
\end{align*}
The slow factor travels at $\delta\omega/q$, tending to
$c_g=\omega'(k_0)$ in the narrowband limit. Thus
\begin{equation}\label{eq:group-speed}
 c_g=\frac{\omega_0}{2k_0}=\frac12c_p.
\end{equation}
Crests move through the group. The group is neither a particular crest nor
a collection of the same fluid particles.

$X=x-c_gt$ follows leading packet translation.
$\partial_t+c_g\partial_x$ measures evolution in that moving frame;
it is the transport derivative in (2.8).

\subsection{Dispersion curvature gives relative phase evolution}
For a nearby longitudinal wavenumber,
\[
 \omega(k_0+q)=\omega_0+c_gq+\tfrac12\omega''(k_0)q^2+\cdots,
 \qquad \omega''(k_0)=-\frac{\omega_0}{4k_0^2}.
\]
The constant is carrier frequency, the linear term gives translation, and
the quadratic term controls relative phase evolution within the packet.
In a PDE it is represented by a second derivative of the envelope.

Transverse curvature has the opposite sign. With dimensionless longitudinal
offset $q$ and transverse offset $r$,
\begin{equation}\label{eq:dispersion-series}
 \frac{\omega(k_0(1+q),k_0r)}{\omega_0}
 =1+\frac q2-\frac{q^2}{8}+\frac{r^2}{4}
 +\frac{q^3}{16}-\frac{3qr^2}{8}+O((|q|+|r|)^4).
\end{equation}
This explains the quadratic terms in NLS and the third derivatives in
Dysthe's equation. To check a PDE coefficient, insert
$e^{i(qx+ry-\nu t)}$ and compare its $\nu$ with this polynomial.

\begin{example}[A scale check]
For wavelength $10\,\mathrm{m}$, $k_0=2\pi/10\,\mathrm{m}^{-1}$.
Taking $g=9.81\,\mathrm{m\,s}^{-2}$ gives
$\omega_0\approx2.48\,\mathrm{s}^{-1}$,
$c_p\approx3.95\,\mathrm{m\,s}^{-1}$, and
$c_g\approx1.98\,\mathrm{m\,s}^{-1}$.
For $a=0.10\,\mathrm{m}$ the steepness is approximately $0.063$.
These are illustrative linear values, not breaking-wave predictions.
\end{example}

\subsection{What the deep-water limit removes}
For a flat bottom at $z=-h$, the no-normal-flow condition is
$\phi_z(-h)=0$. A linear horizontal mode then has vertical factor
$\cosh[k_0(z+h)]$ rather than $e^{k_0z}$. At the surface its logarithmic
vertical derivative is
\[
 \frac{\phi_z}{\phi}=k_0\tanh(k_0h).
\]
The same two linear surface conditions therefore give
\[
 \omega_0^2=gk_0\tanh(k_0h).
\]
For $k_0h\gg1$, $\tanh(k_0h)$ tends to one and
\eqref{eq:carrier-dispersion} is recovered. For $k_0h\ll1$,
$\tanh(k_0h)\approx k_0h$ and $\omega_0\approx k_0\sqrt{gh}$.
The shallow-water limiting phase and group speeds are equal; the deep-water
factor of two is a consequence of a different dispersion law.

This derivation explains why depth is a model assumption, not merely a
descriptive adjective. A change of depth modifies dispersion and the
mean-flow boundary problem, so one cannot retain the deep-water envelope
coefficients unchanged. Dysthe mentions finite-depth extensions in the
introduction, but the equations derived in this paper use infinite depth.

\paperlink{Read (2.7)--(2.8), p.~107, and the dispersive terms of (2.19),
p.~109. Optional reinforcement: \cite{plumb}.}
\checkpoint{Recover the $q^2$ and $r^2$ signs in
\eqref{eq:dispersion-series} from the expansion of $(1+s)^{1/4}$.}
\clearpage
